\section{Variable Length Masked Diffusions: Inference}
\label{sec:FlexMDM_inference}
In this section, we outline inference algorithms for FlexMDM, focusing on two variants: \xxpurple{\textbf{vanilla inference}} and \xxpurple{\textbf{adaptive inference}}. We begin with a brief overview of vanilla and adaptive inference in MDMs.

\textbf{Adaptive inference in MDM.} For the case of MDM, MDM inference proceeds by simulating the rate matrix entries in \ref{eq:rate_mdm}. From a high-level one way this can be done is by (a) independently sampling a subset of masked tokens to unmask and (b) sampling clean tokens from the unmasking posterior. Crucially for what follows, the same guarantee holds for non-independent \emph{adaptive} choices of unmasking indices, e.g., confidence-based: correctness hinges on using the ground-truth unmasking posterior, not on following the rate matrix’s unmasking entries. 
This adaptive inference strategy is widely used due to its empirical performance. We adopt this template and show that FlexMDM inherits the same any-order property.

\textbf{Vanilla inference.} We begin with the \emph{vanilla inference} of FlexMDM, which is obtained by discretizing the CTMC in~\eqref{eq:rate_FlexMDM} using trained neural networks $(f_\theta,g_\theta)$. Choosing an appropriate discretization scheme is crucial, as different schemes can lead to markedly different empirical behavior. We adopt \emph{$\tau$-leaping}—originating in chemical physics and shown to outperform naive Euler discretization for MDMs~\citep{campbell2022continuous}—which batches all events occurring within a fixed interval $[t,t+\tau]$. At a high level, for each discretized step, we simultaneously (Figure~\ref{fig:flexmdm_overview}, right):
\begin{itemize}[leftmargin=*,itemsep=0pt,topsep=0pt]
    \item \coloredhl{xblue}{Unmasking}: For each mask token, sample for every unmasking a number according to the unmasking intensities in the rate matrix. Unmask only if a non-zero entry is returned.
    \item \coloredhl{xxgreen}{Insertion}: Sample the number of \emph{mask-token insertions} from a Poisson distribution parameterized by the insertion rate, then apply those insertions.
\end{itemize}
As the number of steps $\to\infty$, this inference algorithm recovers the CTMC and the discretization error vanishes. Algorithm~\ref{alg:inference} details the full sampler.

\textbf{Adaptive inference.} Notably, one can choose the positions to unmask \coloredhl{xxpurple}{adaptively}. Precisely, at each inference step we select the unmasking positions according to a heuristic rule that prioritizes \coloredhl{xxpurple}{the most confident indices}, where confidence is computed either from the \emph{model’s unmasking posterior} or via a \emph{semi-autoregressive} rule (prioritizing leftmost masks). We find such adaptive choice substantially boosts performance; see Section~\ref{sec:experiment}.

Since unmasking indices in an adaptive no longer trace the transitions described by the rate matrix entries defined in \eqref{eq:rate_FlexMDM}, one might ask whether sampling still guarantees to sample from the target distribution $p_1$ in the infinitesimal limit. The following proposition answers in the affirmative.
\begin{figure}[!t]
\captionsetup{type=algorithm}
\centering
\begin{minipage}[t]{0.55\textwidth}
\vspace{-0.575in}
\begin{algorithm}[H]
\caption*{\textbf{Subroutine 1:} VLMDM inference}
\begin{algorithmic}[1]
\Require Learned functions $(f_\theta, g_\theta)$
\Require Discretization $0 = t_1 < \dots < t_N = 1$
\Require Insertion, Unmasking schedule $\alpha_t,\beta_t$
\State Initialize \(X_{t_1} \gets \varepsilon\)
\For{$k=1$ to $N-1$}
    \State \(\tau \gets t_{k+1} - t_k\)
    \State \xxpurple{\textbf{Invoke Subroutine 2 for unmasking}}
    \For{\xxgreen{\textbf{\(i\) in $[\mathrm{len}(X_{t_k})]$}}}
        \State Set rate $r\gets\tfrac{\dot{\alpha}_{t_k}}{1-\alpha_{t_k}}\cdot\tau$
        \State Sample $\ell\sim\mathrm{Poi}\left(r \cdot g_\theta(X_{t_k}, t_k)[i] \right)$
        \State \xxgreen{\textbf{Insert $\ell$ masks between $X_{t_k}^{i-1}$ and $X_{t_k}^i$}}
    \EndFor
\EndFor
\State \Return $X_{t_N}$
\end{algorithmic}
\end{algorithm}
\end{minipage}\hfill
\begin{minipage}[t]{0.43\textwidth}
\vspace{-0.58in}
\begin{algorithm}[H]
\caption*{\textbf{Subroutine 2: }Unmasking Step}
\begin{algorithmic}[1]
% \Require $X_{t_k},f_\theta,\tau$
\If{\xxpurple{\textbf{vanilla inference}}}:
    \For{$i \in \{i|X_{t_k}^i=\mask\}$ and $v\in\Sigma$}
    \State Set rate $r\gets\frac{\dot{\beta}_{t_k}}{1-\beta_{t_k}}\cdot\tau$
    \State $k_v\sim\mathrm{Poi}(r\cdot f_\theta(X_{t_k}, t_k)[i,v])$
    \If{$^{\exists!}v$ such that $k_v=1$}
    \State Set $X_{t_k}^i \leftarrow v$
    \vspace{-0.02in}
    \EndIf
    \EndFor
\EndIf
\If{\xxpurple{\textbf{adaptive inference}}}:
\State Select $\mathrm{K}$ (the size of $|S|$)
\For{$i \in \{i|X_{t_k}^i=\mask\}$} 
\State Compute confidence $\mathcal{C}^i$
\vspace{-0.025in}
\EndFor
\vspace{-0.03in}
\For{\(i\) in $\mathrm{argmaxK}({C})$}
\State  $X_{t_k}^i \sim \mathrm{Cat}(f_\theta(X_{t_k},t_k)[i])$
\EndFor
\EndIf
\end{algorithmic}
\end{algorithm}
\end{minipage}
\caption{\textbf{VLMDM inference.} At each step we perform \xblue{\textbf{unmasking}} and \xxgreen{\textbf{insertion}}. For \xblue{\textbf{unmasking}}, unmask by $\tau$-leaping (\xxpurple{\textbf{vanilla}}) or by confidence-based selection (\xxpurple{\textbf{adaptive}}). The number of mask tokens to \xxgreen{\textbf{insert}} is drawn from a Poisson distribution. \textbf{Notation}: $\mathrm{Cat}$, $\mathrm{Poi}$ imply the categorical and Poisson distribution, respectively. $\mathrm{argmaxK}(\mathcal{C})$ is the indices set of the $K$ largest components of $\mathcal{C}$. We provide more details in Appendix~\ref{sec:appendix_theory_inference}.} \label{alg:inference}
\vspace{-0.2in}
\end{figure}

\begin{proposition}[Any-order inference, informal]\label{thm:inference}
Consider any sampling scheme that, at each step: (i) unmasks an arbitrary subset of masked positions but draws revealed tokens from the \emph{ground-truth unmasking posterior}; and (ii) applies \emph{insertion} CTMC governed by the ground-truth rate matrix. Then the resulting process samples from the target distribution $p_1$.
\end{proposition}
The formal statement and the proof of Proposition~\ref{thm:inference} are given in Appendix~\ref{sec:appendix_theory_inference}. In words, following the unmasking entries of the rate matrix corresponding to the schedule used in training is \emph{not} necessary to preserve the sampling guarantee. Moreover, the samplers as $N \to \infty$ in Algorithm~\ref{alg:inference} is subsumed by the class in Proposition~\ref{thm:inference}, therefore, assuming access to the ground-truth unmasking posterior and insertion expectation, the corresponding class of algorithms in Algorithm~\ref{alg:inference} samples from $p_1$ up to discretization error.

\textbf{Remark.} A key technical ingredient underlying the rigor of our adaptive inference is that the respective entries of the unmasking posterior of the ground truth rate matrix in Proposition~\ref{thm:inference} do \emph{not} depend on the choice of unmasking schedule $\beta_t$ (the proof is given in Appendix~\ref{sec:appendix_unmasking_posterior}). This independence allows a single model $f_\theta$ to learn all possible unmasking transitions arising along different paths that ultimately connect $p_0$ to $p_1$, thereby \coloredhl{xxpurple}{enabling adaptive unmasking} at inference time. This feature is the same mechanism enabling adaptive inference for MDMs, but for FlexMDMs, proving that it interfaces correctly with insertions is quite subtle. Note that a similar notion—independence of the choice of path—has been introduced in continuous spaces~\cite{albergo2023multi, negrel2025task}. We defer further discussion to Appendix~\ref{sec:appendix_theory_inference}.